\section{Discussion}\label{sec:discussion}

\subsection{Closing the Planning--Execution Gap}

The experiments support the central distinction behind ECR-Cover: planned pointing and effective physical observation are different events. The nominal results show that a planner built around end-to-end confirmation time can change the complete first-coverage trajectory, not only the geometric route. MAG remained an executable reference and supplied a deterministic planning bound, yet ECR-Cover's paired hardware UAC was lower in every nominal block and its T100 tail was shorter. The smaller but resolved UAC difference from calibrated static planning indicates that part of the improvement comes from using a measured execution model and part from the subsequent baseline-preserving route search. These conclusions are empirical; Proposition~\ref{prop:initial} alone guarantees only that the planner did not accept a worse robust model objective.

The repair ablation provides a more direct causal comparison because the initial plan, interruption rule, and confirmation pipeline were held fixed. Repair reduced UAC in nine of ten blocks while the T100-ratio interval remained inside the noninferiority boundary. The adverse block is important: baseline-preserving model acceptance does not eliminate run-to-run execution variation. Its inclusion shows why physical superiority must be evaluated from independent paired runs rather than inferred from the planning rule. At the same time, the pre-recovery witnesses demonstrate that the observed aggregate effect corresponds to actual reassignment. Every repair-on run produced at least one confirmed cell from a migrated task before the paused device returned.

The execution ledger changes what can be claimed from a coverage experiment. A nominal arrival schedule would have supplied a smooth coverage curve but no proof that its steps corresponded to images. Here, each step can be joined to a successful command path, a stability interval, a fresh usable frame, a synchronized measured pose, and a cell projection. This evidence is especially relevant for networked cameras because the control and video paths do not share a single completion signal. ACK, arrival, and frame freshness must be represented separately. The ledger does not make the camera model exact; it makes the relationship between the model and recorded events explicit.

The UAC objective also reveals behavior that makespan alone would hide. ECR-Cover and calibrated static planning had similar nominal T95, yet their UAC and T100 differed. In repair ablation, the main benefit appeared in UAC and T95 while T100 stayed nearly unchanged. This pattern is consistent with moving useful observations earlier without necessarily changing the final interrupted task. The higher load CV shows the accompanying allocation choice: first-observation latency may favor a camera that can acquire a large set early, even when equal task counts would appear more balanced. Because that tradeoff was encoded before the experiment, it is an outcome of the objective rather than a post hoc explanation.

\subsection{Relationship Among Formal and Empirical Evidence}

The results form three evidence layers. The first layer is planning structure. Retaining MAG and the original residual queues makes nonreplacement observable: a time-limited search can always return a complete, already evaluated plan. Equations~\eqref{eq:initial-accept} and \eqref{eq:repair-accept} then give the respective model-relative inequalities. This layer is deterministic once the admissible map and time bounds are fixed.

The second layer is implementation consistency. The state machine and ledger ensure that a coverage write has the fields required by Proposition~\ref{prop:evidence}. The 461-event audit verifies that the frozen implementation followed this rule in all formal runs. Conditional completeness is likewise matched by the observed deployment: the admissible candidates jointly covered all 384 cells, every logical pause ended, and every run reached 384/384 confirmed cells. This observation verifies the premises in the tested setting, not for arbitrary future hardware faults.

The third layer is run-level physical statistics. Paired nominal blocks test whether model-relative planning improvement survives device dynamics and network timing. Repair-on/off blocks test whether the event mechanism improves the realized trajectory. Confidence intervals and sign tests belong to this layer. They neither prove the deterministic propositions nor become consequences of them. Conversely, a significant paired result cannot replace the baseline-preservation proof, because it does not prevent an individual future planning search from accepting a worse modeled plan. The value of ECR-Cover lies in combining the layers without collapsing their meanings.

The four-corner planning core illustrates the need for this separation. Offline sensitivity shows that a $1^\circ$ contraction materially changes candidate set sizes and task counts. It is therefore not a negligible numerical detail. Nevertheless, it remains a particular feasible-map construction. The initial and repair properties would apply to a different common map, while the external accuracy of either map would require independent geometric calibration. Presenting the core as a control-coordinate planning margin, rather than as a physical confidence region, preserves what the data actually support.

\subsection{Limitations}

Eight limitations delimit the present evidence. First, coverage is evaluated in one frozen local geometry. Station positions, mounting rotations, and individual camera intrinsics were not validated against external target ground truth, so model-consistent projection is not the same as absolute world-coordinate accuracy. Second, the $1^\circ$ four-corner intersection is linked to the pose-completion gate and measured telemetry errors; it does not jointly quantify every source of field-of-view uncertainty.

Third, no target trajectory or detection ground truth was available. UAC therefore measures first spatial observation opportunities. It can be interpreted as truncated expected detection delay only under the stationary, uniformly located, certain-detection assumptions stated in Section~\ref{sec:problem}. Fourth, the end-to-end timing margins depend on the relevance of the stratified calibration runs. The 89.6\% independent empirical coverage is consistent with the selected level but does not guarantee performance under nonstationary congestion, altered streaming parameters, or a different device population.

Fifth, the deployment contains four cameras at one site. It does not establish behavior for larger teams, heterogeneous cameras, or geometries with substantially different overlap. Sixth, the formal missions are minute-scale. Long-term thermal drift, encoder bias, camera restart behavior, and multi-day network variability were not tested. Seventh, unavailability was introduced as a controlled logical pause. The protocol tests residual-task reassignment and recovery timing, but it does not cover mechanical damage, corrupted telemetry, partial video degradation, or an unreachable device that never returns.

Eighth, planning was centralized on a soft real-time Windows/Python stack. The measured 200 and 40~ms budgets were met in the frozen experiment, but the data do not establish hard real-time deadlines, distributed consensus behavior, or communication contention in a larger network. These limitations identify the evidence needed for broader deployment claims: external geometric calibration, target-level validation, persistent-fault protocols, larger networks, and longer-duration trials. They are not addressed by expanding the interpretation of the current 70 runs.

